\section{Methods}
The proposed estimator uses penalized least squares. For an observed series
$y_1,\ldots,y_n$, the fitted trend $\hat{\tau}$ is obtained by balancing
fidelity to the observations against a finite-difference roughness penalty. A
typical objective has the form
\[
  \sum_{t=1}^n (y_t-\tau_t)^2 + \lambda \|D^d \tau\|_2^2,
\]
where $D^d$ is a difference operator of order $d$ and $\lambda$ controls
smoothness. This formulation follows the finite-difference penalized smoothing
tradition associated with Whittaker-Henderson graduation, Guerrero's
smoothness index, and modern penalized-spline and perfect-smoother treatments
\citep{whittaker1922graduation,henderson1924graduation,
guerrero2007penalized,eilers1996psplines,eilers2003perfect,wahba1990splines}.

The train-validation selector fits candidate penalized trends on the training
segment, forecasts the validation segment, and selects the candidate with the
lowest validation loss. The time-weighted variant uses a validation loss that
can emphasize later validation observations. The baseline methods are an
HP-style trend \citep{hodrick1997business,king1993lowfrequency}, a
Whittaker-Henderson-style penalized trend
\citep{whittaker1922graduation,henderson1924graduation}, a trailing moving
average, simple exponential smoothing, and a polynomial trend forecaster.
Moving-average and exponential-smoothing benchmarks are included because they
are common forecasting baselines \citep{brown1959forecasting,holt2004exponential,
winters1960forecasting,gardner1985exponential,gardner2006exponential,
hyndman2008exponential}. The polynomial benchmark is a low-dimensional
regression baseline \citep{draper1998regression}.

Guerrero's contribution is used here as a common smoothness scale for the
methods whose fitted trends solve a finite-difference penalized least-squares
problem. The train-validation selector and the time-weighted selector search
over Guerrero smoothness values and choose the value that minimizes their
validation loss. HPTrend and WhittakerTrend do not learn smoothness from the
validation set in this paper; instead, their conventional penalty choices are
converted back to the same Guerrero index so they can be interpreted on the
same model-side scale. Thus the four Guerrero-compatible methods share the
same penalty geometry, but they differ in how the penalty is selected.

The temporal protocol is strict. The series is split into 60 percent training,
20 percent validation, and 20 percent test observations. Hyperparameters are
selected using only the training and validation segments. Final test forecasts
are produced after refitting on the combined training and validation data.
